% Compiles on its own (pulls ../main.tex's preamble via subfiles) or as part of the book.
\documentclass[../main.tex]{subfiles}
\begin{document}

\chapter{How YOLOX Is Built, and Where Fusion Goes}

This is the theory gate for the model. By the end you should be able to draw YOLOX on paper,
label the tensor shapes, and point at the exact place our fusion module attaches.

\section{Three parts}

Every modern one-stage detector has the same three parts.

\begin{enumerate}
  \item \textbf{Backbone} --- eats the image, produces feature maps at three scales.
  \item \textbf{Neck} --- mixes those three scales together so small and large objects are both
        handled well.
  \item \textbf{Head} --- turns feature maps into boxes and class scores.
\end{enumerate}

Our input is $640\times512$. Everything below uses those numbers so the shapes are concrete.

\section{The backbone: CSPDarknet}

The backbone shrinks the image step by step while growing the number of channels. It outputs
three feature maps, named after the stage that produced them:

\begin{center}
\begin{tabular}{llrl}
\toprule
Name & Stride & Size & Holds \\
\midrule
\texttt{dark3} & 8  & $80\times64$, 128 ch & edges, texture, small parts \\
\texttt{dark4} & 16 & $40\times32$, 256 ch & limbs, torsos \\
\texttt{dark5} & 32 & $20\times16$, 512 ch & ``a person-ish thing is here'' \\
\bottomrule
\end{tabular}
\end{center}

\textbf{Stride} is how many input pixels one cell of the feature map covers. Stride 8 means each
cell summarises an $8\times8$ patch. Check the arithmetic: $640/8=80$, $512/8=64$. And
$640/32=20$, $512/32=16$.

Going deeper, you gain \emph{meaning} and lose \emph{location}. \texttt{dark3} knows precisely
where an edge is but not what it belongs to. \texttt{dark5} knows there is a person but only
roughly where.

The backbone is made of four repeated pieces:
\begin{itemize}
  \item \textbf{ConvBNAct} --- convolution, batch-norm, activation. The atom.
  \item \textbf{Bottleneck} --- two ConvBNActs with an optional skip connection.
  \item \textbf{CSPLayer} --- splits the channels in two, sends one half through several
        Bottlenecks, and concatenates. ``CSP'' means Cross-Stage Partial. It gets most of the
        benefit of a deep stack for about half the compute, which is why it is everywhere.
  \item \textbf{SPPBottleneck} --- pools with several kernel sizes ($5,9,13$) and concatenates,
        so one layer sees several receptive-field sizes at once. Sits at the end, in
        \texttt{dark5}.
\end{itemize}

\section{The neck: PAFPN}

A pedestrian 60\,px tall is best detected at one scale; one 300\,px tall at another. The neck
lets the three scales share information. PAFPN does two passes:

\begin{enumerate}
  \item \textbf{Top-down}: upsample \texttt{dark5}, merge into \texttt{dark4}, upsample that,
        merge into \texttt{dark3}. Semantic meaning flows \emph{down} to the high-resolution maps.
  \item \textbf{Bottom-up}: downsample the result and merge back up. Precise location flows
        \emph{up} to the semantic maps.
\end{enumerate}

Out come three maps at the same three strides, but each now informed by the others.

\section{The head: decoupled, anchor-free}

At every cell of every output map, the head answers three questions with three separate small
branches --- this separation is the ``decoupled head'', and it is a large part of why YOLOX
converges faster than earlier YOLOs:

\begin{itemize}
  \item \textbf{class} --- what is it?
  \item \textbf{objectness} --- is there anything here at all?
  \item \textbf{box} --- four numbers giving the box, predicted \emph{directly} (anchor-free:
        no template boxes to hand-tune).
\end{itemize}

Total predictions per image: $80{\cdot}64 + 40{\cdot}32 + 20{\cdot}16 = 5120+1280+320 = 6720$.
Most say ``nothing here.'' Non-maximum suppression then removes duplicates.

\section{SimOTA, and why we do not rewrite it}

During training something must decide \emph{which} of those 6720 cells is responsible for each
real pedestrian. That is \textbf{label assignment}, and YOLOX uses \textbf{SimOTA}: it builds a
cost matrix (classification + box error) between every candidate cell and every ground-truth
object, then dynamically picks the best $k$ cells per object, with $k$ chosen per object rather
than fixed.

We port SimOTA and the losses rather than reimplement them. The reason is methodological, not
lazy: SimOTA is \emph{held fixed across every experiment in this project}. It is not a variable
we ablate. A subtle bug there would corrupt every row of the results table equally and be nearly
impossible to detect, and debugging it would consume weeks that buy no scientific result. We
write the parts we vary; we borrow the parts we hold constant.

\begin{lesson}
There is a real interaction to stay alert to. Label assignment depends on the loss values at each
cell, and fusion changes the features feeding those cells. So a fusion change can shift which
cells become positives, producing an accuracy change that has \emph{nothing to do with better
features}. If an ablation row moves unexpectedly, check the positive/negative split before
believing the architecture story.
\end{lesson}

\section{Where fusion attaches}

Here is the whole point of the chapter. For two modalities we run \emph{two} backbones --- one on
RGB, one on thermal --- and merge before the neck:

\begin{center}
\begin{tabular}{l}
\toprule
RGB $\rightarrow$ backbone$_\text{rgb}$ $\rightarrow$ (\texttt{dark3}, \texttt{dark4}, \texttt{dark5}) \\
thermal $\rightarrow$ backbone$_\text{thm}$ $\rightarrow$ (\texttt{dark3}, \texttt{dark4}, \texttt{dark5}) \\
\midrule
\textbf{fusion} at one or more stages $\rightarrow$ fused (\texttt{dark3}, \texttt{dark4}, \texttt{dark5}) \\
\midrule
$\rightarrow$ PAFPN neck $\rightarrow$ decoupled head $\rightarrow$ boxes \\
\bottomrule
\end{tabular}
\end{center}

Two design facts follow immediately.

\textbf{The fused output must keep the same shapes} the neck already expects
($128/256/512$ channels at strides $8/16/32$). Any fusion operator that changes the channel count
must project back. This is why ``concat then $1\times1$ convolution'' is the workhorse: concat
doubles the channels, and the $1\times1$ halves them back.

\textbf{Which stage you fuse at is Axis A of the ablation}, and by Chapter~1's stride argument it
is also the misalignment story. An 8\,px shift between the modalities is a full cell at
\texttt{dark3} (stride 8), but a quarter of a cell at \texttt{dark5} (stride 32). Fusing deep is
automatically more tolerant of misregistration; fusing shallow demands near-perfect alignment.

Every fusion operator we write must also accept a \emph{missing} thermal stream (the
\texttt{thermal\_valid} flag from the plan amendment), because Soham's composability requirement
says the RGB path must still work when thermal does not arrive.

\section{Build order}

We build bottom-up, testing shapes at each step before adding the next:
\begin{enumerate}
  \item the blocks: ConvBNAct, Bottleneck, CSPLayer, SPPBottleneck;
  \item assemble CSPDarknet, and assert it emits the three shapes in the table above;
  \item the PAFPN neck;
  \item the decoupled head;
  \item ported SimOTA + losses, then a single-stream training run (the RGB-only baseline).
\end{enumerate}
Only then do we go two-stream. A fusion model built on an unverified backbone is impossible to
debug.

\end{document}
